\subsection{Fresh acyclic-support profile rules at nineteen arrows}

The following rules apply to hom-cardinality matrices.  They assert an
obstruction for every composition law with the stated matrix; they do not
replace the finite composition search on the residual matrices.

\begin{lemma}[Proved: recursive zero-second-cohomology rules]
\label{CP2:A-nineteen-profile-H2-rules}
Let $E$ be a full subcategory on at most four objects of a category with
acyclic support.  The following recursive procedure is sufficient to prove
$H^2(BE;\mathbb Q)=0$.  Apply it separately to each connected support
component.  A component passes if one of the following holds:
\begin{enumerate}
\item It is one corner of order at most four.
\item It has $r$ objects and at most $2r+1$ arrows.
\item It has at least two objects and all its nonempty strict hom-sets
are singletons.
\item All its corners have order at most two and its strict support has
height at most one.
\item Its hom-cardinality matrix exhibits an actual initial or terminal
object, including its trivial corner.
\item It has an object $v$ whose strict incoming arrows consist of one
arrow from a trivial corner, or the dual outgoing condition, and the
full deletion of $v$ passes this procedure.
\item It has a singleton pendant whose neighbor's corner has order at
most four, and the full deletion of that pendant passes this procedure.
\end{enumerate}
\end{lemma}
\begin{proof}
The first case is Corollary~\ref{CP2:A-five-types}; the second is
Lemma~\ref{CP2:A-high-H2-cost}.  In the third case the actual support
functor induces an isomorphism in rational second cohomology by
Lemma~\ref{CP2:A-singleton-support}.  An order complex on at most four
vertices has zero second homology: a nonzero simplicial two-cycle needs
at least the boundary of a tetrahedron, whereas a clique complex
containing all four of those triangles also contains the tetrahedron.
The fourth case is Lemma~\ref{CP2:A-semisimple-height}; the fifth is the
usual natural-transformation contraction of a category with an initial
or terminal object.

In the sixth case the full inclusion of the deletion is a homotopy
equivalence by the literal singleton-comma proof of
Lemma~\ref{CP2:A-five18-singleton-deletion}; the proof does not depend on
its historical eighteen-arrow application.  In the seventh case
Lemma~\ref{CP2:A-rational-pendant-deletion} makes that same inclusion a
rational homology equivalence.  Each recursive step removes an object,
so induction terminates.  Disjoint components contribute direct sums
in positive rational cohomology.
\end{proof}

\begin{proposition}[Proved: strict actual reductions at nineteen arrows]
\label{CP2:A-nineteen-actual-profile-reductions}
A nineteen-arrow category of type $\mathbb CP^2$ with acyclic support
cannot satisfy any of the following cardinality conditions:
\begin{enumerate}
\item An extreme object has a nontrivial corner, at least one incident
strict hom-set, and every nonempty incident strict hom-set is a singleton.
\item A trivial singleton pendant at $q$ has another singleton hom-set
in the same direction at $q$, with a different other object.
\item An object has just one strict incoming arrow, from a trivial
corner, or just one strict outgoing arrow, to a trivial corner.
\end{enumerate}
\end{proposition}
\begin{proof}
In the first case all crossing actions are trivial.  The actual corner
retraction in Theorem~\ref{CP2:A-trivial-extreme-corner-retract} preserves
projective-plane type and loses at least one arrow.  The second case is
the actual fold of Proposition~\ref{CP2:A-singleton-pendant-fold}, losing
two arrows.  The third case is the actual singleton-comma deletion,
which loses the deleted identity and the specified strict arrow, hence
at least two arrows.  Each would give a projective-plane category with
at most eighteen arrows, contrary to
Theorem~\ref{CP2:A-global-lower-nineteen}.
\end{proof}

\paragraph{Other uniformly valid predicates.}
The fresh four- and five-object scope generator also applies the support
fundamental-group quotient, the singleton-support comparison, the
semisimple-height bound, Lemmas~\ref{CP2:A-high-two-parallel} and
\ref{CP2:A-high-parallel-cut}, the two-block and fixed-arrow lemmas,
and the small-binary source and sink lemmas.  In a proper block it uses
only the recursive tests of Lemma~\ref{CP2:A-nineteen-profile-H2-rules}.
The tree predicates are exactly
Proposition~\ref{CP2:A-tree-one-large-corner-eleven},
Proposition~\ref{CP2:A-tree-small-binary-corners}, and the actual tree-edge
comparison of Lemma~\ref{CP2:A-five18-tree-edges}; their corner and stalk
bounds are tested literally.  A star with one corner of order at most
eleven and all other corners of order at most two uses the same actual
leaf-orbit reduction to a bounded ordinary act cone.  A trivial pendant
with at least two parallel arrows to a trivial corner is excluded by
Lemma~\ref{CP2:A-trivial-pendant-parallel-circle}.

All-trivial corners give an acyclic category.  Four objects have nerve
dimension at most three.  At five objects, the exact arrow minimum
of Corollary~\ref{A7:exact-object-and-arrow-minima} is thirty-two, so
nineteen arrows cannot realize the projective plane.  The semisimple
height predicate is applied only when strict height is at most three;
the former eighteen-arrow full-five-chain shortcut is not used.

For either object count, the support graph has cycle rank at most six.
Thus the cycle-rank hypotheses of the two parallel-hom lemmas hold
throughout.  Every use of a bounded-homology or rank-one-square lemma is
under the assumed projective-plane type.  The generator's completed
profile count is only a cardinality partition; a residual profile is
never declared excluded without its separate proof or finite gate.
